\scanpage{8}
\mspage{7}

- 13 -
\hfill n° 391

terminer la démonstration, soient $y$ un point de $C$, $z$ un point de $C'=\bigcap_{H\in\uline{\unsure{M}}}D_H(C)$.

Comme $\uline{H}$ est localement fini, le segment fermé d'extrémités $y$, $z$ ne peut rencontrer qu'un nombre fini d'hyperplans de $\uline{H}$ ; il n'y a par suite qu'un nombre fini d'hyperplans $H_1,\ldots,H_q$ de $\uline{H}$ tels que $z\notin D_{H_i}(C)$ pour $1\leq i\leq q$. Par définition de $C'$, les $H_i$ ne sont pas des murs de $C$, donc, par récurrence sur $i$, on voit que $C$ est l'intersection des $D_H(C)$ pour les $H\in\uline{H}$ \uline{distincts des} $H_i$. Cela prouve que $z\in C$, donc $C'\subset C$, et par suite $C'=C$.

COROLLAIRE 2.- \uline{Soit} $\uline{H}$ \uline{un ensemble localement fini d'hyperplans affines.} \uline{Alors tout hyperplan} $H\in\uline{H}$ \uline{est mur mitoyen de deux chambres déterminées par} $\uline{H}$.

Il suffit d'appliquer la prop. 10 à l'ensemble $\uline{H}-\{H\}$.

COROLLAIRE 3.- \uline{Soient} $\uline{H}$ \uline{un ensemble localement fini d'hyperplans affines,} $C$ \uline{une chambre déterminée par} $\uline{H}$, $H_{\unsure{o|0}}$ \uline{un mur de} $C$. \uline{L'intérieur de} $\bar C\cap H_{\unsure{o|0}}$ \uline{dans} $H_{\unsure{o|0}}$ \uline{est une chambre dans} $H_{\unsure{o|0}}$, \uline{formée des points de} $\bar C\cap H_{\unsure{o|0}}$ \uline{qui n'appartiennent à aucun hyperplan} $H\neq H_{\unsure{o|0}}$ \uline{de} $\uline{H}$ ; \uline{ce sont aussi les points} $x$ \uline{de} $\bar C\cap H_{\unsure{o|0}}$ \uline{qui sont adhérents à} $C$ \uline{et à une seule des autres chambres déterminées par} $\uline{H}$, \uline{et cette chambre} $C_1$ \uline{est mitoyenne à} $C$ \uline{et est la même pour tous les points} $x$ \uline{intérieurs à} $\bar C\cap H_{\unsure{o|0}}$.

Le fait que $H_{\unsure{o|0}}$ soit un mur de $C$ entraîne que $C_{\unsure{o|0}}=\bigcap_{H\in\uline{H},\,H\neq H_{\unsure{o|0}}}D_H(C)$ est distincte de $C$ (cor. 1) ; $C$ est donc l'une des deux chambres notées $C'$, $C''$ dans l'énoncé de la prop. 10 ; supposons par exemple que $C=C'$ ; avec les mêmes notations, l'intérieur de $\bar C\cap H_{\unsure{o|0}}$ dans $H_{\unsure{o|0}}$ est $T$, et on a vu dans la prop. 10 que les points de $T$ ne sont adhérents qu'à deux des chambres déterminées par $\uline{H}$, $C'$ et la chambre mitoyenne $C''$. Au contraire, si un point $x\in\bar C\cap H_{\unsure{o|0}}=\bar T$ n'appartient pas à $T$, il appartient à un des hyperplans $H\neq H_{\unsure{o|0}}$ de $\uline{H}$ ; il est
